\scanpage{174}\mspage{29}
% archive sheet number: 174
% Original C page filename: H134-05_134-2_scan174.tex
supposons qu'il y ait exactement deux strates, une strate fermée $X_0$ et $X_1=X-X_0$. D'après le dévissage d'Artin, un faisceau $F$ sur $X$ revient à la donnée d'un faisceau $F_0$ \inserted{$(=i_0^*(F))$} sur $X_0$, d'un faisceau $F_1$ \inserted{$(=i_1^*(F))$} sur $X_1$, et d'un homomorphisme $F_0\longrightarrow i_0^*i_{1*}(F_1)$. Pour que $F$
% Diagram physically in the left margin beside line 7; no enclosing parentheses.
{\reversemarginpar\margin{\diagram{\begin{tikzcd}[ampersand replacement=\&,column sep=1.1em,cells={font=\scriptsize}]
X_0 \arrow[r,hook,"i_0"] \& X \& X_1 \arrow[l,hook']
\end{tikzcd}}}}%
soit $\unsure{\mathcal P|P}$-constructible, il f et s. que $F_0$ et $F_1$ soient loc. constants (plus des conditions de finitude accessoires\ldots), sur $X_0$ et $X_1$ respectivement. Alors (grâce aux hyp. d'équisingularités) il en sera de $\overline m$ de $i_0^*\unsure{i|i'}_{1\unsure{*|,}}(F_1)=\varphi(F_1)$, et \inserted{la cat.} des faisceaux qui nous \unsure{intéressent}\deleted{\unsure{s|t}} s'exprime entièrement en \deleted{\unsure{fonction}} \inserted{termes} de catégories de faisceaux loc. constants sur $X_0$ et $X_1$, \deleted{\uline{\unsure{i.e. de la|i.e. de leurs}} \unsure{d}} \inserted{i.e. des \inserted{seuls} types} d'homotopie de $X_0$ et $X_1$, à celà près qu'il faut encore expliciter la nature du foncteur exact à g. $\varphi$. Je pense que ceci est possible, dans le contexte \inserted{\uline{schématique}} où je me suis placé (techniquement assez sophistiqué), en introduisant un ``voisinage tubulaire étale'' de $X_0$ dans $X$ (qui est un topos fort intéressant, mais pas associé à un schéma). Mais cette construction technique ne fait que para-%
